% GENERATED FILE. Edit canonical lessons, then run npm run generate:books.
% canonical-source-hash: fnv1a64:e115a9b99179aee1
% canonical-lessons: ZH-C123-congming, ZH-C123-renzhen, ZH-C123-keai, ZH-C123-reqing, ZH-C123-guanxin

\chapter{What People Are Like}
\label{ch:zh-what-people-are-like}
\hlchaptermodality{123}

\begin{chapteropening}
\textbf{By the end of this chapter:} \emph{I can say that someone is clever, serious, lovable, warm or caring.}

Complete the last lesson of chapter 123: I can say that someone is clever, serious, lovable, warm or caring.
\end{chapteropening}

\section[cōngming]{cōngming — clever; bright}

\label{lesson:ZH-C123-congming}



\begin{quote}
Before the new one: say the Mandarin for timetable, then the Mandarin for to give a discount.
\end{quote}

\subsection*{You'll want to know: cōngming}
\emph{cōngming} — \textquotedblleft{}clever; bright\textquotedblright{}.

\begin{quote}
\emph{wǒ mèimei hěn cōngming} — My little sister is very bright.
\end{quote}

Its characters come later; for now it is a word for the ear.

\begin{cousinweb}
\emph{cōngming} is built from \emph{cōng} and \emph{míng}, both inherited from Middle Chinese, \emph{tshuwng} and \emph{mjaeng} (Guangyun, in Baxter's transcription). The dictionary glosses the first 'to hear' and 'to discern', and the second 'bright'.
\end{cousinweb}

\begin{grammarlens}[title={What you've built}]
One more word for describing things.
\end{grammarlens}

\subsection*{Your turn}
\begin{itemize}
\raggedright
  \item \emph{Say it:} \emph{cōngming}
  \item \emph{Say it:} \emph{cōngming}, once more
  \item \emph{Say it:} \emph{wǒ mèimei hěn cōngming} — My little sister is very bright.
  \item \emph{Recall:} say the Mandarin for timetable, then the Mandarin for to give a discount, then say \emph{cōngming} again
\end{itemize}

\begin{tcolorbox}[breakable,colback=teal!4,colframe=teal!35!black,title={Before you move on}]
What does cōngming mean? (\textquotedblleft{}Clever; bright\textquotedblright{}.) Say it once more.
\end{tcolorbox}
\section[rènzhēn]{rènzhēn — conscientious; serious}

\label{lesson:ZH-C123-renzhen}



\begin{quote}
Before the new one: say the Mandarin for to give a discount, then the Mandarin for clever.
\end{quote}

\subsection*{You'll want to know: rènzhēn}
\emph{rènzhēn} — \textquotedblleft{}conscientious; serious\textquotedblright{}.

\begin{quote}
\emph{tā gōngzuò hěn rènzhēn} — He takes his work seriously.
\end{quote}

Its characters come later; for now it is a word for the ear.

\begin{cousinweb}
\emph{rènzhēn} is built from \emph{rèn} 'to recognise' and \emph{zhēn} 'true', both inherited from Middle Chinese, \emph{nyinH} and \emph{tsyin} (Guangyun, in Baxter's transcription). How the pair came to mean conscientious is not given there, so it is not settled.
\end{cousinweb}

\begin{grammarlens}[title={What you've built}]
One more word for describing things.
\end{grammarlens}

\subsection*{Your turn}
\begin{itemize}
\raggedright
  \item \emph{Say it:} \emph{rènzhēn}
  \item \emph{Say it:} \emph{rènzhēn}, once more
  \item \emph{Say it:} \emph{tā gōngzuò hěn rènzhēn} — He takes his work seriously.
  \item \emph{Recall:} say the Mandarin for to give a discount, then the Mandarin for clever, then say \emph{rènzhēn} again
\end{itemize}

\begin{tcolorbox}[breakable,colback=teal!4,colframe=teal!35!black,title={Before you move on}]
What does rènzhēn mean? (\textquotedblleft{}Conscientious; serious\textquotedblright{}.) Say it once more.
\end{tcolorbox}
\section[kě'ài]{kě'ài — cute; lovable}

\label{lesson:ZH-C123-keai}



\begin{quote}
Before the new one: say the Mandarin for clever, then the Mandarin for conscientious.
\end{quote}

\subsection*{You'll want to know: kě'ài}
\emph{kě'ài} — \textquotedblleft{}cute; lovable\textquotedblright{}.

\begin{quote}
\emph{zhè zhī māo hěn kě'ài} — This cat is so cute.
\end{quote}

Its characters come later; for now it is a word for the ear.

\begin{cousinweb}
\emph{kě'ài} is built from \emph{kě}, glossed 'to permit', and \emph{ài} 'to love', both inherited from Middle Chinese, \emph{khaX} and \emph{'ojH} (Guangyun, in Baxter's transcription): worth loving.
\end{cousinweb}

\begin{grammarlens}[title={What you've built}]
One more word for describing things.
\end{grammarlens}

\subsection*{Your turn}
\begin{itemize}
\raggedright
  \item \emph{Say it:} \emph{kě'ài}
  \item \emph{Say it:} \emph{kě'ài}, once more
  \item \emph{Say it:} \emph{zhè zhī māo hěn kě'ài} — This cat is so cute.
  \item \emph{Recall:} say the Mandarin for clever, then the Mandarin for conscientious, then say \emph{kě'ài} again
\end{itemize}

\begin{tcolorbox}[breakable,colback=teal!4,colframe=teal!35!black,title={Before you move on}]
What does kě'ài mean? (\textquotedblleft{}Cute; lovable\textquotedblright{}.) Say it once more.
\end{tcolorbox}
\section[rèqíng]{rèqíng — warm; welcoming}

\label{lesson:ZH-C123-reqing}



\begin{quote}
Before the new one: say the Mandarin for conscientious, then the Mandarin for lovable.
\end{quote}

\subsection*{You'll want to know: rèqíng}
\emph{rèqíng} — \textquotedblleft{}warm; welcoming\textquotedblright{}.

\begin{quote}
\emph{wǒ de lǎoshī hěn rèqíng} — My teacher is very warm.
\end{quote}

Its characters come later; for now it is a word for the ear.

\begin{cousinweb}
\emph{rèqíng} is built from \emph{rè} 'hot' and \emph{qíng} 'feeling', both inherited from Middle Chinese, \emph{nyet} and \emph{dzjeng} (Guangyun, in Baxter's transcription); the dictionary explains the second, after an older book, as the breath in a person that has desires.
\end{cousinweb}

\begin{grammarlens}[title={What you've built}]
One more word for describing things.
\end{grammarlens}

\subsection*{Your turn}
\begin{itemize}
\raggedright
  \item \emph{Say it:} \emph{rèqíng}
  \item \emph{Say it:} \emph{rèqíng}, once more
  \item \emph{Say it:} \emph{wǒ de lǎoshī hěn rèqíng} — My teacher is very warm.
  \item \emph{Recall:} say the Mandarin for conscientious, then the Mandarin for lovable, then say \emph{rèqíng} again
\end{itemize}

\begin{tcolorbox}[breakable,colback=teal!4,colframe=teal!35!black,title={Before you move on}]
What does rèqíng mean? (\textquotedblleft{}Warm; welcoming\textquotedblright{}.) Say it once more.
\end{tcolorbox}
\section[guānxīn]{guānxīn — to care about}

\label{lesson:ZH-C123-guanxin}



\begin{quote}
Before the new one: say the Mandarin for lovable, then the Mandarin for warm-hearted.
\end{quote}

\subsection*{You'll want to know: guānxīn}
\emph{guānxīn} — \textquotedblleft{}to care about\textquotedblright{}.

\begin{quote}
\emph{māma hěn guānxīn wǒ} — Mum really cares about me.
\end{quote}

Its characters come later; for now it is a word for the ear.

\begin{cousinweb}
\emph{guānxīn} is built from \emph{guān}, the word for shutting, and \emph{xīn} 'heart', both inherited from Middle Chinese, \emph{kwaen} and \emph{sim} (Guangyun, in Baxter's transcription). The dictionary glosses the first as the bar that shuts a door; how the pair came to mean caring is not settled.
\end{cousinweb}

\begin{grammarlens}[title={What you've built}]
One more word for describing things.
\end{grammarlens}

\subsection*{Your turn}
\begin{itemize}
\raggedright
  \item \emph{Say it:} \emph{guānxīn}
  \item \emph{Say it:} \emph{guānxīn}, once more
  \item \emph{Say it:} \emph{māma hěn guānxīn wǒ} — Mum really cares about me.
  \item \emph{Recall:} say the Mandarin for lovable, then the Mandarin for warm-hearted, then say \emph{guānxīn} again
\end{itemize}

\begin{tcolorbox}[breakable,colback=teal!4,colframe=teal!35!black,title={Before you move on}]
What does guānxīn mean? (\textquotedblleft{}To care about\textquotedblright{}.) Say it once more.
\end{tcolorbox}
